% ============================================================================
% PHẦN 2 (2022 - ĐẦU 2023): CĂN CHỈNH MÔ HÌNH BẰNG HỌC TĂNG CƯỜNG (RLHF & PPO)
% ============================================================================
\sectionintro{2}{Alignment với RLHF và PPO}{2022--đầu 2023}{Sở thích cặp, Reward Model, KL regularization và kiến trúc bốn mạng.}

\subsection{2.1. Giới thiệu Học tăng cường (RL) \& Reward Model}

% ----------------------------------------------------------------------------
% SLIDE: TỔNG QUAN HỌC TĂNG CƯỜNG (RL)
% ----------------------------------------------------------------------------
\begin{frame}{Tổng quan về Học tăng cường (RL)}
  \lead{Khung làm việc tương tác giữa tác tử và môi trường (Sutton \& Barto, 2018).}
  \subhead{Các thành phần cốt lõi của RL trong mô hình ngôn ngữ}
  \small
  \begin{itemize}\setlength{\itemsep}{5pt}
    \item \textbf{Agent (Tác tử):} Mô hình ngôn ngữ cần tối ưu hóa chính sách $\pi_\theta$.
    \item \textbf{Environment (Môi trường):} Ngữ cảnh hội thoại nhận phản hồi.
    \item \textbf{State $s_t$ \& Action $a_t$:} State là prompt và chuỗi token đã sinh $(x, y_{<t})$; Action là token tiếp theo $y_t \sim \pi_\theta(\cdot \mid s_t)$.
    \item \textbf{Reward (Phần thưởng) $r$:} Điểm số đánh giá chất lượng toàn bộ phản hồi.
  \end{itemize}
  \vspace{0.15cm}
  \textbf{Mục tiêu:} Tối đa hóa kỳ vọng phần thưởng tích lũy $\E_{\tau \sim \pi_\theta} [R(\tau)]$.
  \par\vspace{0pt}
\end{frame}

% ----------------------------------------------------------------------------
% SLIDE: TẠI SAO DÙNG RL SAU SFT
% ----------------------------------------------------------------------------
\begin{frame}{Tại sao dùng RL sau SFT?}
  \lead{Dùng tín hiệu thưởng để chấm điểm và định hướng chất lượng câu trả lời.}
  \subhead{Động lực áp dụng RL sau khi SFT}
  \small
  \begin{itemize}\setlength{\itemsep}{5pt}
    \item \textbf{Vượt qua giới hạn Teacher Forcing:}\\
      SFT chỉ học bắt chước từng token độc lập, không đánh giá được tổng thể câu trả lời.
    \item \textbf{Đánh giá chuỗi tự sinh hoàn chỉnh:}\\
      RL cho phép mô hình tự sinh trọn vẹn câu trả lời $\hat{y} \sim \pi_\theta$, sau đó nhận điểm thưởng.
    \item \textbf{Tối ưu hóa trực tiếp theo chất lượng:}\\
      Thưởng điểm cao ($r > 0$) cho phản hồi hữu ích, phạt điểm thấp ($r < 0$) đối với phản hồi độc hại hoặc sai lệch.
  \end{itemize}
  \par\vspace{0pt}
\end{frame}

% ----------------------------------------------------------------------------
% SLIDE: DỮ LIỆU ĐỐI CHIẾU CẶP
% ----------------------------------------------------------------------------
\begin{frame}{Thu thập dữ liệu cặp đối chiếu}
  \lead{InstructGPT chuyển từ gán nhãn đơn lẻ sang quan hệ thứ tự cặp (Ouyang et al., 2022).}
  \subhead{Quy trình xây dựng tập $\mathcal{D}_{\text{pair}}$}
  \small
  \begin{itemize}\setlength{\itemsep}{6pt}
    \item \textbf{Lấy mẫu cặp phản hồi:} Truy vấn $x \sim \mathcal{D}_{\text{prompt}}$, sinh cặp đầu ra $(y_1, y_2) \sim \pi_{\text{SFT}}(\cdot \mid x)$.
    \item \textbf{Gán nhãn ưu tiên (Pairwise Comparison):} Chuyên gia chọn $y_w \succ y_l$.
      \begin{itemize}\small\setlength{\itemsep}{2pt}
        \item $y_w$ (Preferred / Mẫu dương): Phản hồi hữu ích, an toàn hơn.
        \item $y_l$ (Dispreferred / Mẫu âm): Phản hồi lỗi, độc hại hoặc lạc đề.
      \end{itemize}
  \end{itemize}
  \par\vspace{0pt}
\end{frame}

% ----------------------------------------------------------------------------
% SLIDE: TẠI SAO SO SÁNH CẶP
% ----------------------------------------------------------------------------
\begin{frame}{Tại sao sử dụng so sánh cặp?}
  \lead{Loại bỏ phương sai cá nhân và trôi ngưỡng chấm điểm tuyệt đối.}
  \subhead{Ưu thế thống kê của quan hệ nhị phân}
  \small
  \begin{itemize}\setlength{\itemsep}{6pt}
    \item \textbf{Hạn chế của thang điểm tuyệt đối ($1\text{--}10$):}\\
      Phương sai cao giữa các chuyên gia; ngưỡng chấm bị trôi dạt (Calibration Drift) theo thời gian.
    \item \textbf{Tính nhất quán của so sánh nhị phân ($y_w \succ y_l$):}\\
      Độ đồng thuận thống kê vượt trội; dễ chuẩn hóa trên quy mô dữ liệu lớn.
    \item \textbf{Bổ sung Mẫu âm $y_l$ (Negative Sample):}\\
      Cung cấp ranh giới đối chiếu rõ ràng để triệt tiêu hành vi sai lệch.
  \end{itemize}
  \par\vspace{0pt}
\end{frame}

% ----------------------------------------------------------------------------
% SLIDE: TỔNG QUAN REWARD MODEL
% ----------------------------------------------------------------------------
\begin{frame}{Mô hình Phần thưởng (Reward Model)}
  \lead{Thước đo tự động hóa sở thích và tiêu chuẩn đánh giá của con người.}
  \subhead{Cấu trúc và cơ chế hoạt động}
  \small
  \begin{itemize}\setlength{\itemsep}{6pt}
    \item \textbf{Đầu vào (Input):} Cặp gồm câu truy vấn $x$ và câu trả lời hoàn chỉnh $y$.
    \item \textbf{Đầu ra (Output):} Điểm số vô hướng $r_\phi(x, y) \in \mathbb{R}$ phản ánh độ chính xác và mức ưa thích.
    \item \textbf{Vai trò cốt lõi:}\\
      Đóng vai trò hàm mục tiêu tự động (Automated Proxy), thay thế con người liên tục chấm điểm cho hàng triệu phản hồi trong quá trình huấn luyện RL.
  \end{itemize}
  \par\vspace{0pt}
\end{frame}

% ----------------------------------------------------------------------------
% SLIDE: SƠ ĐỒ HUẤN LUYỆN REWARD MODEL
% ----------------------------------------------------------------------------
\begin{frame}{Huấn luyện Reward Model}
  \lead{Học điểm vô hướng $r_\phi(x,y)\in\mathbb{R}$ từ cặp mẫu dương và mẫu âm.}
  \vspace{0.15cm}
  \begin{center}
    \resizebox{0.96\textwidth}{!}{%
      \begin{tikzpicture}
        \node[data, text width=2.8cm, minimum height=1.2cm] (prompt) at (0, 0) {
          \textbf{Truy vấn $x$}\\
          (Từ tập $\mathcal{D}_{\text{pair}}$)
        };

        \node[data, fill=softgreen, draw=tealgreen, text width=3.3cm, minimum height=1.1cm] (yw) at (4.2, 1.15) {
          \textbf{Mẫu dương $y_w$}\\
          (Positive Sample)
        };

        \node[data, fill=softred, draw=stanfordred, text width=3.3cm, minimum height=1.1cm] (yl) at (4.2, -1.15) {
          \textbf{Mẫu âm $y_l$}\\
          (Negative Sample)
        };

        \node[trainmodel, text width=3.8cm, minimum height=3.1cm] (rm) at (9.0, 0) {
          \textbf{Reward Model $r_\phi$}\\
          (Reward Model)\\[6pt]
          \normalfont\scriptsize Khởi tạo từ $\pi_{\text{SFT}}$;\\
          Thay đầu từ vựng $|\mathcal{V}|$ bằng\\
          \textbf{Đầu Chiếu Vô hướng}\\
          $W_r \in \mathbb{R}^{d \times 1}$ tại \texttt{<EOS>}
        };

        \node[rewardbox, text width=2.9cm, minimum height=1.0cm] (rw) at (13.3, 1.15) {
          Điểm mẫu dương\\$\mathbf{r_\phi(x, y_w) \in \mathbb{R}}$
        };

        \node[lossbox, fill=softred, draw=stanfordred, text width=2.9cm, minimum height=1.0cm] (rl) at (13.3, -1.15) {
          Điểm mẫu âm\\$\mathbf{r_\phi(x, y_l) \in \mathbb{R}}$
        };

        \node[lossbox, text width=3.6cm, minimum height=1.6cm] (btloss) at (17.5, 0) {
          \textbf{Bradley--Terry Loss}\\[3pt]
          $\mathcal{L}_{\text{RM}} = -\log \sigma(r_w - r_l)$\\[3pt]
          \normalfont\scriptsize Cực đại hóa hiệu số $r_w - r_l$
        };

        \coordinate (bus_p) at (1.9, 0);
        \draw[line width=1.05pt, draw=navyblue!85!black] (prompt.east) -- (bus_p);
        \draw[flow] (bus_p) |- (yw.west);
        \draw[flow] (bus_p) |- (yl.west);

        \draw[flow] (yw.east) -- ([yshift=1.15cm]rm.west);
        \draw[flow] (yl.east) -- ([yshift=-1.15cm]rm.west);

        \draw[flow] ([yshift=1.15cm]rm.east) -- (rw.west);
        \draw[flow] ([yshift=-1.15cm]rm.east) -- (rl.west);

        \draw[flow] (rw.east) -| (btloss.north);
        \draw[flow] (rl.east) -| (btloss.south);

        \draw[gradflow] (btloss.south) -- (17.5, -2.35) -- node[below, font=\scriptsize\rmfamily\bfseries, text=stanfordred] {Backpropagation $\nabla_\phi \mathcal{L}_{\text{RM}}$ cập nhật tham số $\phi$} (9.0, -2.35) -- (rm.south);
      \end{tikzpicture}%
    }
  \end{center}
  \vspace{-0.05cm}
  \small
  \textbf{Nguyên lý:} Cặp $(x, y_w)$ và $(x, y_l)$ cùng đi qua $r_\phi$; gradient kéo tăng $r_\phi(x, y_w)$ ($\uparrow$) và hạ thấp $r_\phi(x, y_l)$ ($\downarrow$).
  \par\vspace{0pt}
\end{frame}

% ----------------------------------------------------------------------------
% SLIDE: MÔ HÌNH BRADLEY-TERRY
% ----------------------------------------------------------------------------
\begin{frame}{Mô hình Bradley--Terry và Loss}
  \lead{Maximum Likelihood Estimation cho quan hệ thứ tự cặp.}
  \begin{block}{Bradley--Terry (1952) và Loss $\mathcal{L}_{\text{RM}}(\phi)$}
    \small
    \[
      \fitmath{%
      P(y_w \succ y_l \mid x) = \sigma\big(r_\phi(x, y_w) - r_\phi(x, y_l)\big), \qquad
      \mathcal{L}_{\text{RM}}(\phi) = -\E_{(x, y_w, y_l)} \Big[ \log \sigma\big(r_\phi(x, y_w) - r_\phi(x, y_l)\big) \Big]
      }
    \]
  \end{block}

  \subhead{Ý nghĩa thực tiễn}
  \small
  \begin{itemize}\setlength{\itemsep}{5pt}
    \item \textbf{Chuyển đổi thành phân loại nhị phân:}\\
      Mô hình hóa xác suất ưu tiên thông qua hàm sigmoid của hiệu số điểm $r_\phi(x, y_w) - r_\phi(x, y_l)$.
    \item \textbf{Nới rộng khoảng cách chất lượng:}\\
      Cực đại hóa xác suất giúp nới rộng chênh lệch điểm giữa phản hồi tốt $r_\phi(x, y_w)$ và phản hồi kém $r_\phi(x, y_l)$.
  \end{itemize}
  \par\vspace{0pt}
\end{frame}

\subsection{2.2. Định luật Goodhart, KL regularization (\texorpdfstring{$\beta \DKL$}{β D\textunderscore KL}) \& Chu trình RL}

% ----------------------------------------------------------------------------
% SLIDE: ĐỊNH LUẬT GOODHART & REWARD HACKING
% ----------------------------------------------------------------------------
\begin{frame}{Định luật Goodhart và Reward Hacking}
  \lead{Tối đa hóa cực đoan một hàm định giá xấp xỉ sẽ làm suy biến phân phối (Gao et al., 2023).}
  \subhead{Rủi ro khi tối ưu trực tiếp $r_\phi(x, y)$}
  \small
  \begin{itemize}\setlength{\itemsep}{5pt}
    \item \textbf{Định luật Goodhart:} Thước đo trở thành mục tiêu tối ưu sẽ mất tính tin cậy.
    \item \textbf{Sai số ngoài phân phối (OOD Error):}\\
      $r_\phi$ chỉ tin cậy gần $\pi_{\text{SFT}}$; ra xa vùng huấn luyện, $r_\phi$ cho điểm cao giả tạo.
    \item \textbf{Hiện tượng Reward Hacking:}\\
      Policy sinh văn bản vô nghĩa hoặc lặp từ nhưng vẫn đạt $r_\phi$ cực đại.
  \end{itemize}
  \vspace{0.20cm}
  {\small\bfseries\color{reforange}$\implies$ Cần giữ phân phối output model không đi quá xa khỏi model sau SFT.\par}
  \par\vspace{0pt}
\end{frame}

% ----------------------------------------------------------------------------
% SLIDE: ĐIỀU CHUẨN VÙNG TIN CẬY BẰNG KL DIVERGENCE
% ----------------------------------------------------------------------------
\begin{frame}{Điều chuẩn vùng tin cậy bằng KL Divergence}
  \lead{Neo giữ Policy trong vùng tin cậy của mô hình tham chiếu $\pi_{\text{ref}}$.}
  \begin{block}{Khoảng cách KL phân phối có điều kiện}
    \[
      \fitmath{%
      \DKL\big(\pi_\theta(\cdot \mid x) \parallel \pi_{\text{ref}}(\cdot \mid x)\big) = \E_{y \sim \pi_\theta} \left[ \log \frac{\pi_\theta(y \mid x)}{\pi_{\text{ref}}(y \mid x)} \right]
      }
    \]
  \end{block}

  \subhead{Ý nghĩa của khoảng cách $\DKL$}
  \small
  \begin{itemize}\setlength{\itemsep}{4pt}
    \item $\DKL \to 0$: Phân phối của Policy trùng khớp hoàn toàn với mô hình sau SFT.
    \item $\DKL \to \infty$: Phân phối Policy trôi dạt quá xa, có nguy cơ suy biến do Reward Hacking.
  \end{itemize}
  \vspace{0.15cm}
  \scriptsize\color{refgray}
  \textbf{Chú thích:}\\
  $\bullet$ \texttt{model\_ref} ($\pi_{\text{ref}}$): Là mô hình giữ nguyên trọng số sau pha SFT.\\
  $\bullet$ \texttt{model\_phi} ($\pi_\theta$): Là mô hình đang được huấn luyện.
  \par\vspace{0pt}
\end{frame}

% ----------------------------------------------------------------------------
% SLIDE: RLHF OBJECTIVE
% ----------------------------------------------------------------------------
\begin{frame}{RLHF Objective}
  \lead{Tối đa hóa phần thưởng kết hợp KL regularization (Ouyang et al., 2022).}
  \begin{block}{Mục tiêu RLHF có điều chuẩn KL}
    \[
      \fitmath{%
      \max_{\pi_\theta} \; J_{\text{RLHF}}(\pi_\theta) = \E_{x \sim \mathcal{D},\; y \sim \pi_\theta} \Big[ r_\phi(x, y) - \beta \log \frac{\pi_\theta(y \mid x)}{\pi_{\text{ref}}(y \mid x)} \Big]
      }
    \]
  \end{block}
  \par\vspace{0pt}
\end{frame}

% ----------------------------------------------------------------------------
% SLIDE: SƠ ĐỒ Online RL Loop
% ----------------------------------------------------------------------------
\begin{frame}{Online RL Loop}
  \lead{Policy tự sinh mẫu $\hat y\sim\pi_\theta$ rồi nhận phản hồi để cập nhật.}
  \vspace{0.12cm}
  \begin{center}
    \resizebox{0.96\textwidth}{!}{%
      \begin{tikzpicture}
        \node[data, text width=3.1cm, minimum height=1.3cm] (state) at (0, 0) {
          \textbf{Trạng thái $s_t$}\\
          Truy vấn $x$ \& tiền tố\\tự sinh $\hat{y}_{<t}$
        };

        \node[trainmodel, text width=3.5cm, minimum height=1.3cm] (policy) at (4.6, 0) {
          \textbf{Policy $\pi_\theta(\hat{y}_t \mid s_t)$}\\
          Mô hình Ngôn ngữ\\đang tối ưu hóa
        };

        \node[data, fill=softgreen, draw=tealgreen, text width=3.4cm, minimum height=1.3cm] (action) at (9.3, 0) {
          \textbf{Quỹ đạo tự sinh}\\
          Rút mẫu $\hat{y} \sim \pi_\theta(\cdot \mid x)$\\
          (On-Policy Rollout)
        };

        \node[rewardbox, text width=3.5cm, minimum height=1.3cm] (env) at (14.0, 0) {
          \textbf{Định giá Phần thưởng}\\
          $R_{\text{total}} = r_\phi(x, \hat{y}) - \beta \DKL$
        };

        \node[lossbox, text width=4.8cm, minimum height=1.35cm] (update) at (12.5, -2.15) {
          \textbf{Advantage Function $A(x, \hat{y})$}\\[2pt]
          $A(x, \hat{y}) = R_{\text{total}} - b(x)$\\[2pt]
          \normalfont\tiny $A > 0$: Tăng xác suất ($\uparrow$) \ | \ $A < 0$: Giảm xác suất ($\downarrow$)
        };

        \draw[flow] (state.east) -- (policy.west);
        \draw[flow] (policy.east) -- (action.west);
        \draw[flow] (action.east) -- (env.west);

        \draw[flow] (env.south) -- (14.0, -1.15) -| ([xshift=0.8cm]update.north);
        \draw[gradflow] (update.west) -- node[midway, above=3pt, font=\scriptsize\rmfamily\bfseries, text=stanfordred] {Cập nhật Gradient $\nabla_\theta J(\theta)$} (4.6, -2.15) -- (policy.south);
      \end{tikzpicture}%
    }
  \end{center}
  \vspace{0.05cm}
  \small
  \textbf{Khắc phục Exposure Bias:} Policy được tối ưu trực tiếp trên trạng thái tự sinh $\hat{s}_t \sim \pi_\theta$ và nhận phạt âm ($A < 0$) ngay khi đi chệch quỹ đạo.
  \par\vspace{0pt}
\end{frame}

\subsection{2.3. Policy Gradient, Thuật toán PPO (Value Network \texorpdfstring{$V_\psi$}{V\textunderscore ψ}) \& Hệ thống 4 Mô hình}

% ----------------------------------------------------------------------------
% SLIDE: POLICY GRADIENT CHO SINH CHUỖI NGÔN NGỮ
% ----------------------------------------------------------------------------
\begin{frame}{Policy Gradient cho sinh chuỗi ngôn ngữ}
  \lead{REINFORCE giải quyết bài toán không vi phân qua phép lấy mẫu rời rạc (Williams, 1992).}
  \begin{block}{Đồng nhất thức Log-derivative Trick}
    \[
      \fitmath{%
      \nabla_\theta J(\theta) = \nabla_\theta \sum_{y} \pi_\theta(y \mid x) R(x, y) = \E_{y \sim \pi_\theta} \Big[ R(x, y) \, \nabla_\theta \log \pi_\theta(y \mid x) \Big]
      }
    \]
  \end{block}

  \subhead{Bản chất toán học}
  \small
  \begin{itemize}\setlength{\itemsep}{5pt}
    \item Lấy mẫu rời rạc $y \sim \pi_\theta$ không thể truyền gradient trực tiếp: $\frac{\partial y}{\partial \theta}$ không tồn tại.
    \item Áp dụng $\nabla_\theta \pi_\theta = \pi_\theta \nabla_\theta \log \pi_\theta$, đạo hàm kỳ vọng trở thành kỳ vọng của gradient log-xác suất.
    \item Cho phép ước lượng Monte Carlo không chệch qua sinh mẫu (Rollout).
  \end{itemize}
  \par\vspace{0pt}
\end{frame}

% ----------------------------------------------------------------------------
% SLIDE: ADVANTAGE VÀ VAI TRÒ CỦA BASELINE
% ----------------------------------------------------------------------------
\begin{frame}{Advantage và vai trò của Baseline}
  \lead{Trừ hàm baseline $b(x)$ để giảm phương sai mà không gây chệch kỳ vọng.}
  \begin{block}{Policy Gradient với Advantage $A(x, y)$}
    \[
      \fitmath{%
      \nabla_\theta J(\theta) = \E_{y \sim \pi_\theta} \Big[ \underbrace{\big(R_{\text{total}}(x, y) - b(x)\big)}_{\text{Advantage } A(x, y)} \, \nabla_\theta \log \pi_\theta(y \mid x) \Big]
      }
    \]
  \end{block}

  \subhead{Đặc tính của hàm Baseline $b(x)$}
  \small
  \begin{itemize}\setlength{\itemsep}{5pt}
    \item \textbf{Bảo toàn kỳ vọng:} $\E_{y}[b(x) \nabla_\theta \log \pi_\theta] = b(x)\nabla_\theta(1) = 0$.
    \item \textbf{Phân định hướng cập nhật:}
      \begin{itemize}\small\setlength{\itemsep}{2pt}
        \item $A(x, y) > 0$: Quỹ đạo tốt hơn kỳ vọng $\longrightarrow$ Tăng xác suất ($\uparrow$).
        \item $A(x, y) < 0$: Quỹ đạo kém hơn kỳ vọng $\longrightarrow$ Giảm xác suất ($\downarrow$).
      \end{itemize}
  \end{itemize}
  \par\vspace{0pt}
\end{frame}

% ----------------------------------------------------------------------------
% SLIDE: PPO-CLIP: KIỂM SOÁT BƯỚC CẬP NHẬT
% ----------------------------------------------------------------------------
\begin{frame}{Thuật toán PPO: cơ chế Clipping}
  \lead{Giới hạn tỷ số xác suất để cập nhật ổn định trong vùng tin cậy (Schulman et al., 2017).}
  \begin{block}{Hàm mục tiêu PPO-Clip}
    \[
      \fitmath{%
      \mathcal{J}_{\text{PPO}}(\theta) = \E_{t} \Big[ \min \Big( \rho_t(\theta) \hat{A}_t, \; \clip\big(\rho_t(\theta), 1 - \epsilon, 1 + \epsilon\big) \hat{A}_t \Big) \Big]
      }
    \]
    với tỷ số xác suất: $\rho_t(\theta) = \dfrac{\pi_\theta(y_t \mid x, y_{<t})}{\pi_{\theta_{\text{old}}}(y_t \mid x, y_{<t})}$.
  \end{block}

  \subhead{Cơ chế chặn gradient}
  \small
  \begin{itemize}\setlength{\itemsep}{5pt}
    \item $\rho_t(\theta)$: Tỷ số đo lường mức thay đổi xác suất giữa chính sách mới và cũ.
    \item Khi $\rho_t(\theta)$ vượt ngoài $[1-\epsilon, 1+\epsilon]$ ($\epsilon = 0.2$), gradient bị triệt tiêu về $0$.
    \item Ngăn ngừa các bước cập nhật quá lớn làm sụp đổ chính sách.
  \end{itemize}
  \par\vspace{0pt}
\end{frame}

% ----------------------------------------------------------------------------
% SLIDE: MẠNG CRITIC TRONG PPO
% ----------------------------------------------------------------------------
\begin{frame}{Mạng Critic trong PPO}
  \lead{Ước lượng baseline động tại từng vị trí token để tính Advantage.}
  \begin{block}{Hàm mất mát của Value Network $V_\psi$}
    \[
      \fitmath{%
      \hat{A}_t = R_t - V_\psi(s_t), \qquad \mathcal{L}_{\text{Critic}}(\psi) = \E_{t} \Big[ \big(R_t - V_\psi(s_t)\big)^2 \Big]
      }
    \]
  \end{block}

  \subhead{Vai trò và rào cản tài nguyên}
  \small
  \begin{itemize}\setlength{\itemsep}{5pt}
    \item \textbf{Token-level Baseline:} Dự báo phần thưởng kỳ vọng từ trạng thái $s_t$, giảm mạnh phương sai của $\hat{A}_t$.
    \item \textbf{Gánh nặng tham số:} Mạng $V_\psi$ có cùng quy mô với $\pi_\theta$ (chiếm gấp đôi VRAM GPU).
    \item \textbf{Sai số xấp xỉ:} $V_\psi$ dự báo lệch khiến tín hiệu $\hat{A}_t$ méo mó, gây cập nhật sai hướng.
  \end{itemize}
  \par\vspace{0pt}
\end{frame}

% ----------------------------------------------------------------------------
% SLIDE: SƠ ĐỒ 4 MÔ HÌNH PPO
% ----------------------------------------------------------------------------
\begin{frame}{RLHF + PPO: bốn mô hình phối hợp}
  \lead{Actor, Reference, Reward Model và Critic cùng hoạt động trên GPU.}
  \vspace{0.08cm}
  \begin{center}
    \resizebox{0.96\textwidth}{!}{%
      \begin{tikzpicture}
        \node[data, text width=2.4cm, minimum height=1.2cm] (prompt) at (0, 0) {
          \textbf{Truy vấn $x$}\\
          (Từ tập dữ liệu)
        };

        \node[trainmodel, text width=3.5cm, minimum height=1.35cm] (actor) at (4.2, 1.35) {
          \textbf{1. Actor / Policy $\pi_\theta$}\\
          (Đang huấn luyện)\\
          \normalfont\scriptsize Tự sinh chuỗi $y \sim \pi_\theta(\cdot \mid x)$
        };

        \node[frozenmodel, text width=3.5cm, minimum height=1.35cm] (ref) at (4.2, -1.35) {
          \textbf{2. Reference $\pi_{\text{ref}}$}\\
          (Cố định tham số)\\
          \normalfont\scriptsize Tính $\log \pi_{\text{ref}}(y \mid x)$
        };

        \node[data, fill=softgreen, draw=tealgreen, text width=2.5cm, minimum height=1.2cm] (resp) at (8.5, 0) {
          \textbf{Chuỗi tự sinh $y$}\\
          (On-policy Rollout)
        };

        \node[frozenmodel, text width=3.6cm, minimum height=1.35cm] (rm) at (12.8, 1.35) {
          \textbf{3. Reward Model $r_\phi$}\\
          (Cố định tham số)\\
          \normalfont\scriptsize Định giá $r_\phi(x, y)$
        };

        \node[trainmodel, fill=softyellow, draw=amberorange!90!black, text width=3.6cm, minimum height=1.35cm] (critic) at (12.8, -1.35) {
          \textbf{4. Critic / Value $V_\psi$}\\
          (Đang huấn luyện)\\
          \normalfont\scriptsize Ước lượng nền $V_\psi(s_t)$
        };

        \node[lossbox, text width=3.7cm, minimum height=2.2cm] (ppo) at (17.5, 0) {
          \textbf{Tối ưu hóa PPO}\\[3pt]
          $R_{\text{total}} = r_\phi - \beta \DKL$\\[2pt]
          $\hat{A}_t = R_{\text{total}} - V_\psi(s_t)$\\[3pt]
          \normalfont\scriptsize Cắt tỉa tỷ số $\clip(\rho_t)$
        };

        \coordinate (bus_p) at (1.8, 0);
        \draw[line width=1.05pt, draw=navyblue!85!black] (prompt.east) -- (bus_p);
        \draw[flow] (bus_p) |- (actor.west);
        \draw[flow] (bus_p) |- (ref.west);

        \draw[flow] (actor.east) -| (resp.north);
        \draw[flow] (resp.south) |- (ref.east);

        \coordinate (bus_r) at (10.3, 0);
        \draw[line width=1.05pt, draw=navyblue!85!black] (resp.east) -- (bus_r);
        \draw[flow] (bus_r) |- (rm.west);
        \draw[flow] (bus_r) |- (critic.west);

        \draw[flow] (rm.east) -| ([xshift=-0.4cm]ppo.north);
        \draw[flow] (critic.east) -| ([xshift=-0.4cm]ppo.south);
        \draw[flow] (ref.south) -- (4.2, -2.45) -- node[below, font=\scriptsize\rmfamily, text=slate] {Truyền $\log \pi_{\text{ref}}(y \mid x)$ để tính phân kỳ $\beta \DKL(\pi_\theta \parallel \pi_{\text{ref}})$} (17.5, -2.45) -- (ppo.south);

        \draw[gradflow] (ppo.north) -- (17.5, 2.45) -- node[above, font=\scriptsize\rmfamily\bfseries, text=stanfordred] {Cập nhật Gradient PPO đồng thời cho Policy $\pi_\theta$ và Critic $V_\psi$} (4.2, 2.45) -- (actor.north);
      \end{tikzpicture}%
    }
  \end{center}
  \par\vspace{0pt}
\end{frame}

\subsection{2.4. Tổng kết Phần 2: Hạn chế của RLHF + PPO}

% ----------------------------------------------------------------------------
% SLIDE: HẠN CHẾ CỦA RLHF + PPO
% ----------------------------------------------------------------------------
\begin{frame}{Hạn chế của RLHF + PPO}
  \begin{enumerate}\setlength{\itemsep}{4pt}
    \item \textbf{4-Model Memory Bottleneck:}\\
      \small Duy trì đồng thời bốn mạng ($\pi_\theta, \pi_{\text{ref}}, r_\phi, V_\psi$) trên GPU, chiếm gấp đôi VRAM.
    \item \textbf{Autoregressive Rollout Bottleneck:}\\
      \small Sinh mẫu tự hồi quy từng token On-policy chiếm $>70\%$ thời gian tính toán.
    \item \textbf{Variance and Approximation Error:}\\
      \small Sai số xấp xỉ từ $r_\phi$ và Critic $V_\psi$ khiến gradient có phương sai lớn và bất ổn định.
  \end{enumerate}
  \par\vspace{0pt}
\end{frame}
